\documentclass[12pt,a4paper]{article}

% ---------- Fonts: Times New Roman family ----------
\usepackage[T1]{fontenc}
\usepackage[utf8]{inputenc}
\usepackage{newtxtext,newtxmath}   % Times-compatible text + math
\usepackage{textcomp}

% ---------- Page geometry ----------
\usepackage[a4paper,top=1in,bottom=1in,left=1.1in,right=1.1in]{geometry}
\usepackage{setspace}
\onehalfspacing
\usepackage{microtype}                 % better justification, fewer overfull boxes
\setlength{\emergencystretch}{3em}      % mop up residual tight lines

% ---------- Structure / layout ----------
\usepackage{titlesec}
\usepackage{titling}
\usepackage{booktabs}
\usepackage{longtable}
\usepackage{array}
\usepackage{tabularx}
\usepackage{ragged2e}
\usepackage{enumitem}
\usepackage{xcolor}
\usepackage{fancyhdr}
\usepackage{parskip}
\usepackage[hidelinks]{hyperref}
\hypersetup{colorlinks=true,linkcolor=black,urlcolor=blue!55!black,citecolor=black}

% ---------- Heading sizes: main 16pt, sub 14pt, subsub 12pt ----------
\titleformat{\section}{\normalfont\fontsize{16}{19}\bfseries}{\thesection}{0.6em}{}
\titleformat{\subsection}{\normalfont\fontsize{14}{17}\bfseries}{\thesubsection}{0.6em}{}
\titleformat{\subsubsection}{\normalfont\fontsize{12}{15}\bfseries}{\thesubsubsection}{0.6em}{}
\titlespacing*{\section}{0pt}{16pt}{8pt}
\titlespacing*{\subsection}{0pt}{12pt}{6pt}
\titlespacing*{\subsubsection}{0pt}{10pt}{4pt}

% ---------- Header / footer ----------
\pagestyle{fancy}
\fancyhf{}
\renewcommand{\headrulewidth}{0.4pt}
\renewcommand{\footrulewidth}{0.2pt}
\fancyhead[L]{\small AR933 --- Claims Requiring Reconstruction}
\fancyhead[R]{\small Why / How / What}
\fancyfoot[L]{\small Confidential --- Inventor's Suggestions}
\fancyfoot[R]{\small Page \thepage\ of \pageref{LastPage}}
\usepackage{lastpage}

% ---------- Helpers ----------
\usepackage{url}
\DeclareUrlCommand\code{\urlstyle{tt}}
\makeatletter
\g@addto@macro\UrlBreaks{\do\_\do\:\do\-\do\,\do\/\do\.}
\makeatother
\newcommand{\clm}[1]{\textbf{Claim~#1}}
\definecolor{hdrgray}{gray}{0.92}
\newcolumntype{L}[1]{>{\RaggedRight\arraybackslash}p{#1}}

% Callout boxes
\usepackage{tcolorbox}
\newtcolorbox{advisory}[1][]{colback=red!4,colframe=red!45!black,
  boxrule=0.6pt,arc=2pt,left=6pt,right=6pt,top=5pt,bottom=5pt,#1}
\newtcolorbox{notebox}[1][]{colback=blue!3,colframe=blue!35!black,
  boxrule=0.6pt,arc=2pt,left=6pt,right=6pt,top=5pt,bottom=5pt,#1}
% Box for proposed replacement claim language (the "What")
\newtcolorbox{whatbox}[1][]{colback=green!4,colframe=green!40!black,
  boxrule=0.6pt,arc=2pt,left=6pt,right=6pt,top=5pt,bottom=5pt,#1}

\begin{document}

% =====================================================================
% TITLE PAGE
% =====================================================================
\begin{titlepage}
\centering
\vspace*{1.2cm}
{\fontsize{13}{16}\selectfont CLAIMS REQUIRING RECONSTRUCTION\par}
\vspace{0.4cm}
{\color{gray}\rule{0.75\textwidth}{0.8pt}}\par
\vspace{0.9cm}
{\fontsize{20}{24}\bfseries Which Claims to Re-Draft, and\\[4pt]
the Exact Wording to Use\par}
\vspace{0.6cm}
{\fontsize{14}{18}\selectfont \itshape ``A System and Method for Privacy-Preserving\\[2pt]
Edge-Based Threat Detection in a Domestic Environment''\par}
\vspace{0.9cm}
{\color{gray}\rule{0.75\textwidth}{0.8pt}}\par
\vspace{1.4cm}

\begin{tabular}{@{}r >{\raggedright\arraybackslash}p{0.62\textwidth}@{}}
\textbf{Application reference} & AR933 (Complete Specification, Form~2)\\[3pt]
\textbf{Claims to reconstruct} & Claims 1, 5, and 7 only\\[3pt]
\textbf{Claims requiring no change} & Claims 2, 3, 4, 6, 8, 9, 10\\[3pt]
\textbf{Inventors} & Gonugondla Veera Manikanta; Dr Swathi Tangi\\[3pt]
\textbf{Prepared for} & Adv.\ Pranav Bhat (IN/PA~4580), Patent Agent\\[3pt]
\textbf{Prepared by} & Gonugondla Veera Manikanta (Inventor)\\[3pt]
\textbf{Source of record} & Repository \code{Manikanta25055/Garuda_26}, branch \code{main}\\[2pt]
 & \href{https://github.com/Manikanta25055/Garuda_26}{\code{https://github.com/Manikanta25055/Garuda_26}}\\[3pt]
\textbf{Document status} & Confidential --- for prosecution use\\
\end{tabular}

\vfill
{\small This document suggests re-drafting \textbf{only} the claims that the working
implementation does not support as presently worded. For each such claim it states
\textbf{Why} the current wording is not enabled, \textbf{How} to fix it, and \textbf{What}
exact replacement language could be used. All other claims are already aligned, or will be
brought into alignment by code changes the inventor is making, and need no change. These are
suggestions from the inventor's side and are open to discussion.\par}
\vspace{0.8cm}
\end{titlepage}

% =====================================================================
\section{How to Read This Document}
% =====================================================================

\begin{notebox}
\textbf{These are the inventor's suggestions, not instructions.} The reconstructions proposed
here reflect what the working code can honestly support today. They are offered to help align
the claims with the implementation --- they are not a direction to the agent. If any claim
should be kept as originally drafted for strategy, scope, or prosecution reasons, please flag
it and we will discuss before anything is changed. Where a divergence can be closed by a code
change instead of a claim amendment, the inventor will make that change.
\end{notebox}

Three of the ten claims recite behaviour that the actual working code does not perform,
and cannot be made to perform by drafting alone or by a small code change. Those three ---
\clm{1}, \clm{5}, and \clm{7} --- must be reconstructed. This document gives, for each:

\begin{itemize}[leftmargin=1.4em,itemsep=2pt,topsep=2pt]
  \item \textbf{Why} --- the specific clause that is not enabled, and the technical reason
        the implementation cannot deliver it as worded.
  \item \textbf{How} --- the drafting approach: what to re-scope, narrow, or replace.
  \item \textbf{What} --- ready-to-insert replacement language, shown in a shaded box.
\end{itemize}

Every other claim is either already implemented exactly as written, or differs only in a
parameter value that the inventor is correcting in code so the claim stays as drafted.
Those claims need no change from our side. Section~\ref{sec:scope} lists the full
disposition so the scope is clear.

% =====================================================================
\section{Scope --- Which Claims to Touch}\label{sec:scope}
% =====================================================================

\begin{center}
\renewcommand{\arraystretch}{1.25}
\begin{longtable}{@{}c L{2.7cm} L{7.2cm}@{}}
\toprule
\textbf{Claim} & \textbf{Action} & \textbf{Reason}\\
\midrule
\endfirsthead
\toprule
\textbf{Claim} & \textbf{Action} & \textbf{Reason}\\
\midrule
\endhead
\bottomrule
\endfoot
\textbf{1} & \textbf{Reconstruct} & Two clauses not enabled: the secondary network's recited function, and the absolute restriction on video egress. See Sec.~\ref{sec:c1}.\\
\textbf{2} & Keep (verify data) & Wording stands. The 50~fps figure depends on re-measurement --- see the advisory in Sec.~\ref{sec:c5}.\\
\textbf{3} & Keep & Code corrected to the 3-consecutive-frame gate; claim now matches.\\
\textbf{4} & Keep & Code corrected to the 0.60 / 0.25--0.35 asymmetric thresholds; claim now matches.\\
\textbf{5} & \textbf{Reconstruct} & The ``less than 4 milliseconds'' figure is physically unattainable for the real secondary models. See Sec.~\ref{sec:c5}.\\
\textbf{6} & Keep & Power-draw limit measured and satisfied.\\
\textbf{7} & \textbf{Reconstruct} & Independent method claim; mirrors \clm{1} and carries the same two defects. See Sec.~\ref{sec:c7}.\\
\textbf{8} & Keep & Thermal limit satisfied.\\
\textbf{9} & Keep & 20~ms hardware inference latency measured and satisfied.\\
\textbf{10} & Keep & Operating-mode flags under a single lock implemented as claimed.\\
\end{longtable}
\end{center}

\begin{notebox}
\textbf{One-line summary:} we suggest re-drafting \clm{1}, \clm{5}, and \clm{7} using the
shaded ``What'' text below, and leaving every other claim untouched --- subject to your review.
\end{notebox}

% =====================================================================
\section{Claim 1 --- System Claim (Reconstruct)}\label{sec:c1}
% =====================================================================

\clm{1} carries two clauses that the code does not support. They are independent of each
other; both must be re-drafted.

% ---------------------------------------------------------------
\subsection{Clause 1A --- The Secondary Network's Recited Function}
% ---------------------------------------------------------------

\subsubsection{Why (present wording is not enabled)}
The claim recites the secondary network as one that re-checks the primary's own class label:

\begin{quote}\itshape
``\ldots at least one secondary verification neural network configured to independently
re-evaluate a cropped image region \ldots and to generate a verification output
\textbf{confirming or rejecting a classification} \textbf{assigned by the primary
detection neural network}.''
\end{quote}

The implemented secondary network does not do this. It is a binary \emph{Safe-versus-Weapon}
classifier (MobileNetV2) combined with a depth-variance anti-spoof (MiDaS Small), and it is
run on \emph{person} crops only. Weapon-class detections from the primary are treated as
final outputs of the primary and are never routed to the secondary
(\code{garuda_cascade.py:439--447}). The secondary therefore never ``confirms or rejects''
the primary's weapon-class label; it answers a different question --- ``is this person
associated with a weapon, and is the person a real three-dimensional subject rather than a
photo or screen.'' As worded, the clause describes a function the system does not perform and
is not enabled for the weapon classes.

\subsubsection{How (drafting approach)}
Re-scope the secondary's recited function to what it actually produces: a binary
safe-versus-threat classification of the cropped region \emph{plus} a depth-variance
(anti-spoof) measure, with the alert gated on that combined output. Delete the language that
implies re-verification of the primary's own class label.

\subsubsection{What (replacement language)}
\begin{whatbox}
``\ldots at least one secondary verification neural network 109 configured to receive a
cropped image region corresponding to the detected region of interest and to generate a
verification output, the verification output comprising (i) a binary safe-versus-threat
classification of the cropped image region and (ii) a depth-variance measure indicative of
whether the cropped image region corresponds to a three-dimensional physical subject or to a
substantially planar surface, the host processor 103 being configured to gate the alert
signal on the said verification output;''
\end{whatbox}

% ---------------------------------------------------------------
\subsection{Clause 1B --- The Restriction on Video Egress}
% ---------------------------------------------------------------

\subsubsection{Why (present wording is not enabled)}
The closing \emph{wherein} of \clm{1} recites an absolute bar on sending frames off the box:

\begin{quote}\itshape
``\ldots to restrict transmission of the sequence of video frames captured by the camera
sensor to \textbf{any device external} to the host processor and the neural processing
unit, thereby providing the privacy-preserving edge-based threat detection.''
\end{quote}

The system streams live MJPEG / WebRTC video to remote browsers
(\code{Garuda_web.py:90, 966--968}); raw frames are obscured only in \emph{privacy} mode. In
normal operation the frame sequence does leave the host, so an absolute ``any device
external'' restriction is contradicted by a shipped feature and is not enabled.

\subsubsection{How (drafting approach)}
Narrow the restriction from ``any external device'' to the boundary that is truly held:
frames and inference outputs never reach a \emph{third-party} server; they stay within the
user-owned trust boundary (evidence goes by SSH/SFTP to a user-owned store, no cloud). This
is the honest and defensible scope. (If instead the applicant prefers the absolute wording,
it can be kept only by a code change that forces the privacy blur on \emph{all} streamed
frames in every mode --- the inventor can implement this; note it to the agent.)

\subsubsection{What (replacement language)}
\begin{whatbox}
``\ldots and to restrict transmission of the sequence of video frames captured by the camera
sensor 101, and of any inference output derived therefrom, to any server outside a
user-owned trust boundary, such that neither the video frames nor the inference outputs are
transmitted to a third-party server, thereby providing the privacy-preserving edge-based
threat detection.''
\end{whatbox}

% =====================================================================
\section{Claim 5 --- Secondary-Stage Latency (Reconstruct)}\label{sec:c5}
% =====================================================================

\subsubsection{Why (present wording is not enabled)}
\clm{5} recites:

\begin{quote}\itshape
``\ldots executes the at least one secondary verification neural network within a
\textbf{combined latency of less than 4 milliseconds} per cropped image region \ldots''
\end{quote}

The two real secondary models measure MobileNetV2 $\approx 36$~ms and MiDaS~Small
$\approx 29$~ms per crop on the Hailo-8L (\code{garuda_cascade.py:451, 457}). A combined
secondary latency under 4~ms is not attainable for these models on this accelerator --- it is
a model-size / hardware limit, not something a code change can reach. The figure is
unsupported and must be re-drafted, not coded around.

\subsubsection{How (drafting approach)}
Two truthful options; pick one:
\begin{itemize}[leftmargin=1.4em,itemsep=3pt,topsep=2pt]
  \item \textbf{Option~1 --- state the real number.} Replace ``less than 4 milliseconds''
        with the actual measured combined secondary latency (of order 30--40~ms), re-measured
        with the live secondary in the loop, and soften the surrounding ``real-time'' wording
        accordingly.
  \item \textbf{Option~2 --- claim the fast path.} Recite the lightweight RGB-variance
        anti-spoof proxy, which genuinely runs in under 1~ms without invoking the NPU
        (\code{garuda_cascade.py:88--90, 129--143}), as the real-time verification stage, and
        describe the MobileNet/MiDaS path as a higher-assurance mode. This keeps a truthful
        sub-millisecond figure.
\end{itemize}

\subsubsection{What (replacement language)}
\begin{whatbox}
\textbf{Option 1 (measured value):} ``\ldots executes the at least one secondary verification
neural network 109 within a combined latency of not more than \textbf{[MEASURED]}
milliseconds per cropped image region, the combined latency being measured independently of a
frame processing rate of the primary detection neural network 107 \ldots''
\\[6pt]
\textbf{Option 2 (fast-path proxy):} ``\ldots executes, on the said at least a third one of
the plurality of processing cores and without invoking the neural processing unit 105, a
variance-based anti-spoof verification stage within a latency of less than 1 millisecond per
cropped image region, the said verification stage rejecting a cropped image region whose
pixel-intensity variance is below a spoof-rejection threshold \ldots''
\end{whatbox}

\begin{advisory}
\textbf{Data note affecting Claims 2 and 5.} In the production server the secondary worker is
presently a placeholder that logs and returns without running a network
(\code{Garuda_web.py:445--478}). Any headline figure captured on that build --- the
$\approx 52$~fps ``under full concurrent load'' (\clm{2}) and the secondary-stage latency
(\clm{5}) --- was recorded with the second network \emph{not executing}. Before either number
is relied on for prosecution, re-measure with a live secondary in the loop. This is a
measurement action, not a drafting action; the \clm{2} \emph{wording} can stand if the number
holds.
\end{advisory}

% =====================================================================
\section{Claim 7 --- Method Claim (Reconstruct)}\label{sec:c7}
% =====================================================================

\clm{7} is the independent method claim and mirrors \clm{1}. It carries the same two defects,
in step form, and must receive the same two reconstructions.

\subsubsection{Why (present wording is not enabled)}
Two steps repeat the \clm{1} problems:
\begin{itemize}[leftmargin=1.4em,itemsep=2pt,topsep=2pt]
  \item ``\emph{independently re-evaluating \ldots confirming or rejecting a classification
        assigned by the primary detection neural network}'' --- same issue as Clause~1A.
  \item ``\emph{restricting \ldots transmission of the sequence of video frames \ldots to any
        device external to the host processor and the neural processing unit}'' --- same issue
        as Clause~1B.
\end{itemize}

\subsubsection{How (drafting approach)}
Apply the identical re-scoping used for \clm{1}, cast as method steps.

\subsubsection{What (replacement language)}
\begin{whatbox}
\textbf{Verification step (replaces the ``re-evaluating\ldots confirming or rejecting''
step):}\\
``generating, by the at least one secondary verification neural network 109, a verification
output from the cropped image region, the verification output comprising (i) a binary
safe-versus-threat classification of the cropped image region and (ii) a depth-variance
measure indicative of whether the cropped image region corresponds to a three-dimensional
physical subject or to a substantially planar surface, and gating the alert signal on the
said verification output;''
\\[6pt]
\textbf{Egress step (replaces the ``restricting\ldots to any device external'' step):}\\
``restricting, by the host processor 103, transmission of the sequence of video frames
captured by the camera sensor 101, and of any inference output derived therefrom, to any
server outside a user-owned trust boundary, such that neither the video frames nor the
inference outputs are transmitted to a third-party server, thereby providing the
privacy-preserving edge-based threat detection.''
\end{whatbox}

% =====================================================================
\section{Closing}
% =====================================================================

Our suggestion is to reconstruct \clm{1}, \clm{5}, and \clm{7} using the shaded ``What'' text
above, and to leave \clm{2}, \clm{3}, \clm{4}, \clm{6}, \clm{8}, \clm{9}, and \clm{10} as
drafted --- the inventor is correcting the code so those claims remain true. The
only non-drafting item is the \clm{2}/\clm{5} re-measurement flagged in the advisory above.
These are proposals from the inventor's side; if any claim should be preserved as originally
drafted, we would like to discuss it with you first. All file-and-line references are against
\code{Manikanta25055/Garuda_26}, branch \code{main}
(\href{https://github.com/Manikanta25055/Garuda_26}{\code{https://github.com/Manikanta25055/Garuda_26}}),
and can be verified independently.

\vspace{0.8cm}
\noindent\rule{\textwidth}{0.4pt}
\vspace{0.2cm}

\noindent{\small Prepared by the inventor from a line-by-line source audit. Suggestions only
--- open to discussion. Confidential.}

\end{document}
